\documentclass[class=llncs, crop=false]{standalone}

\input{preamble.tex}
\input{macros.tex}

\begin{document}
\section{Eunoia}\label{sec:eunoia}
% fig:eo-syntax is in intro.tex so it floats to the top of page 2
% ---------------------------------------------------------
This section presents the fragment of Eunoia targeted by our
translation, together with the elaboration step that resolves
$n$-ary applications to normal form.
%
Eunoia currently has no published formal semantics:
its meaning is defined operationally by Ethos.
%
The presentation here was informed by the
Ethos user manual~\cite{EthosUser_manualmdMain},
the Ethos source, and through discussions with the
\cvc{5}/Eunoia developers.
%
% We accordingly do not give a typing relation or
% algorithm for Eunoia, nor a model-theoretic semantics
% analogous to the one in SMT-LIB~2.6~\cite{Barrett2015-standard}.


% ---------------------------------------------------------
\subsection{Syntax, Commands, and Signatures}\label{sec:eo-sig}
% ---------------------------------------------------------
%
\Autoref{fig:eo-syntax} defines the concrete syntax for Eunoia,
where the metavariable $s$ ranges over a set of \emph{symbols}.
%
A term is either a symbol $s$, a literal $ℓ$
(integer, rational, or string), an \emph{application}
$\eapp{s}{\vec t}$, or a \emph{binder application}
$\eapp{s}{\epar{\vec ρ}\ t}$.
%
A \emph{parameter} $\epar{s\ t\ \maybe{ν}}$ declares a symbol
$s$ of type $t$ with an optional \emph{parameter attribute} $ν$.
Similarly, a symbol may have a \emph{constant attribute} $α$
that specifies how its $n$-ary applications are elaborated
to binary form.
%
Eunoia provides a fixed set of \emph{built-in symbols}
$β$ that operate on the syntactic structure of terms;
their role is explained alongside the constructs that
use them (rule declarations in \autoref{sec:eo-rules}
and elaboration in \autoref{sec:eo-elab}).
%
Commands are divided into \emph{standard commands}
and \emph{proof commands}.

% ------------------------------------------------------
\begin{table}[t!]
	\centering
	\caption{Declarations produced by each Eunoia command.
		Parts in $\maybe{\cdot}$ are present only when the
		corresponding attribute is supplied. Each program case
		$c_i = \ecase{l_i}{r_i}$ contributes one rewrite rule.}
	\label{tab:cmd-decls}
	\small
	\begin{tabular}{@{}>{$}p{0.56\linewidth}<{$}@{\ }>{$}p{0.40\linewidth}<{$}@{}}
		\hline
		\textnormal{\textbf{Command}} & \textnormal{\textbf{Declarations}} \\
		\hline
		\epar{\kw{declare-const}\ s\ t\ \maybe{α}}
		                              & \{s : t\}\,\maybe{∪\,\{s ∷ α\}} \\
		\epar{\kw{declare-parameterized-const}\ s\ \epar{\vec ρ}\ t\ \maybe{α}}
		                              & \{s\prn{\vec ρ} : t\}\,\maybe{∪\,\{s ∷ α\}} \\
		\epar{\kw{declare-consts}\ ℓ\ t}
		                              & \textnormal{binds the literal class~$ℓ$ to type~$t$} \\
		\epar{\kw{define}\ s\ \epar{\vec ρ}\ t\ \maybe{\attr{type}\ t'}}
		                              & \{s\prn{\vec ρ} ↪ t\}\,\maybe{∪\,\{s\prn{\vec ρ} : t'\}} \\
		\epar{\kw{program}\ s\ \epar{\vec ρ}\ \attr{signature}\ \epar{\vec t}\ t'\ \epar{c_1 \ldots c_m}}
		                              & \{s\prn{\vec ρ} : \earr{\vec t\ t'}\} ∪ \{l_i ↪ r_i \mid c_i = \ecase{l_i}{r_i}\} \\
		\epar{\kw{declare-rule}\ s\ \epar{\vec ρ}\ \ldots}
		                              & \{s\prn{\vec ρ} : τ^★\}\textnormal{\ (\autoref{sec:eo-rules})} \\
		\epar{\kw{include}\ μ}
		                              & \textnormal{declarations of module~$μ$} \\
		\epar{\kw{assume}\ s\ φ}\textnormal{\,/\,}\epar{\kw{assume-push}\ s\ φ}
		                              & \{s : \eprf{φ}\} \\
		\epar{\kw{step}\ s\ φ\ \attr{rule}\ r\ \maybe{\attr{args}\ \epar{\vec t}}\ \maybe{\attr{premises}\ \epar{\vec p}}}
		                              & \{s : \eprf{φ},\ s ↪ \eapp{r}{\vec t\ \vec p}\} \\
		\hline
	\end{tabular}
\end{table}
%
A \emph{signature} $Σ$ is a set of \emph{declarations},
where each declaration is either a \emph{typing} $s\prn{\vec ρ} : t$,
an \emph{attribution} $s\prn{\vec ρ} ∷ α$,
or a \emph{rewrite rule} $ℓ ↪ r$ with parameters $\vec \rho$.
%
\autoref{tab:cmd-decls} maps each command to
the set of declarations it contributes.
%
Any sequence of standard commands $\vec δ$ thus
determines a signature: the union of the
declarations contributed by each command.
%
A \emph{proof script} $Υ = \prn{\vec δ ⨾ \vec π}$
pairs a sequence of standard commands $\vec δ$
(the \emph{preamble}) with a sequence of proof
commands $\vec π$ (the \emph{body}); $Υ$ likewise
determines a signature, which we call a \emph{proof}.
%

\subsection{Rule Declarations and Proofs}\label{sec:eo-rules}
%
Two of the core built-in symbols play a role specific
to rule declarations.
%
The guard $\eoreqs{t_1}{t_2}{t_3}$ evaluates to $t_3$
only when $t_1$ and $t_2$ are syntactically identical,
and is otherwise stuck; it lets a rule make its
conclusion conditional on a side-condition.
%
The marker $\eoquote{t}$ flags $t$ as a \emph{quoted
	argument} whose value may appear in the return type
of the rule, giving Eunoia's mechanism for
object-level dependent types.
%
%
Given an inference rule declaration
$\epar{\kw{declare-rule}\ s\ \epar{\vec ρ}\ \ldots\ \conclusion φ}$,
the derivation of its contributed type declaration
$s\prn{\vec\rho} : \tau^\star$ is given thus:
%
\begin{enumerate}
	\item Given
	      $\requires{\ecase{l_1}{r_1}\ldots\ecase{l_k}{r_k}}$,
	      define $φ^★$ below. Otherwise, $φ^★ ≔ φ$.
	      $$φ^★ ≔ \eoreqs{l_1}{r_1}{
			      \epar{
				      \,\ldots\,\eoreqs{l_k}{r_k}{φ}}
		      }
	      $$

	\item If the declaration provides no premises,
	      let $τ^★ ≔ \eprf{φ^★}$.
	      Otherwise, define $τ^★$ thus:
	      \begin{enumerate}
		      \item Given
		            $\attr{premises}\ \epar{\plur ψ m}$,
		            let:
		            $$τ^★ ≔ \earr{\eprf{ψ_1}\,\ldots\,\eprf{ψ_m}\ \eprf{φ^★}}$$
		            %
		      \item Given $\premiselist ψ f$ for
		            some associative symbol $f$, redefine:
		            $$φ^★ ≔\eoreqs{\eapp{\eo{is\_list}}{f\ ψ}}{\ev{true}}{φ^★}$$
		            and let $τ^★ ≔ \earr{\eprf{ψ}\ \eprf{φ^★}}$.
	      \end{enumerate}

	\item If the declaration provides
	      $\attr{args}\ \epar{\plur t n}$,
	      then the type contributed for $s$ is:
	      $$s\prn{\vec ρ}
		      : \earr{\eoquote{t_1}\,\ldots\,\eoquote{t_n}\ τ^★}
	      $$
	      Otherwise, $s\prn{\vec ρ}: τ^★$.
\end{enumerate}

\begin{example}\label{ex:eo-rules}
	Consider four rule declarations from
	\texttt{cpc/rules/Booleans.eo}.
	%
	\begin{eobox}
		\lstinputlisting[style=eo, aboveskip=0pt, belowskip=0pt]{listings/rule-modus-ponens.eo}%
		\listsep
		\lstinputlisting[style=eo, aboveskip=0pt, belowskip=0pt]{listings/rule-cnf-implies-pos.eo}%
		\listsep
		\lstinputlisting[style=eo, aboveskip=0pt, belowskip=0pt]{listings/rule-reordering.eo}%
		\listsep
		\lstinputlisting[style=eo, aboveskip=0pt, belowskip=0pt]{listings/rule-and-intro.eo}%
	\end{eobox}
	%
	\noindent
	The $\ev{modus\_ponens}$ rule has only $\attr{premises}$ and
	$\attr{conclusion}$, so item~2a gives:
	$$
		\ev{modus\_ponens}(\vec ρ)  :
		{\color{eoPunct}{\texttt{(}}}\ev{->}\
		\eprf{\ev{F1}}\
		\eprf{\eapp{\ev{=>}}{\ev{F1}\ \ev{F2}}}\
		\eprf{\ev{F2}}
		{\color{eoPunct}{\texttt{)}}}
	$$
	%
	The $\ev{cnf\_implies\_pos}$ rule has
	$\attr{args}\ \epar{\eapp{\ev{=>}}{\ev{F1}\ \ev{F2}}}$;
	item~3 makes this an
	$\eo{quote}$-typed parameter that pattern-matches at
	the point of use, where $φ$ is the term given by
	its $\attr{conclusion}$ attribute:
	$$
		\ev{cnf\_implies\_pos}(\vec ρ)  :
		{\color{eoPunct}{\texttt{(}}}\ev{->}\
		\eoquote{\eapp{\ev{=>}}{\ev{F1}\ \ev{F2}}}\
		\eprf{φ}
		{\color{eoPunct}{\texttt{)}}}
	$$
	%
	The $\ev{reordering}$ rule combines $\attr{args}$
	with $\attr{requires}$: item~1 builds the guarded
	conclusion
	$φ^★ ≔ \eoreqs{\eapp{\eo{list\_minclude}}{\ev{or}\ \ev{C1}\ \ev{C2}}}{\ev{true}}{\ev{C2}}$,
	item~2a adds the premise $\ev{C1}$, and item~3 wraps
	the $\eo{quote}$-typed argument $\ev{C2}$:
	$$
		\ev{reordering}(\vec ρ)  :
		{\color{eoPunct}{\texttt{(}}}\ev{->}\
		\eoquote{\ev{C2}}\
		\eprf{\ev{C1}}\
		\eprf{φ^★}
		{\color{eoPunct}{\texttt{)}}}
	$$
	%
	The $\ev{and\_intro}$ rule has
	$\premiselist{\ev F}{\ev{and}}$, so item~2b guards
	the conclusion with
	$φ^★ ≔ \eoreqs{\eapp{\eo{is\_list}}{\ev{and}\ \ev F}}{\ev{true}}{\ev F}$
	and takes the premise list as a single $\eprf{\ev F}$
	argument:
	$$
		\ev{and\_intro}(\vec ρ)  :
		{\color{eoPunct}{\texttt{(}}}\ev{->}\
		\eprf{\ev F}\
		\eprf{φ^★}
		{\color{eoPunct}{\texttt{)}}}
	$$
\end{example}


% ---------------------------------------------------------
\subsection{Elaboration}\label{sec:eo-elab}
% ---------------------------------------------------------
%
Unlike standard LF-style frameworks, Eunoia supports
$n$-ary applications of binary symbols, with the
unrolling determined by a \emph{constant attribute}
$α$ (\autoref{fig:eo-syntax}) assigned to the symbol
at its declaration.
%
After elaboration, every application of a declared
constant $s$ uses the built-in higher-order
application operator $\mtt{\_}$, so that
$\eapp{s}{t_1\ t_2}$ becomes $\eho{\eho{s}{t_1}}{t_2}$;
only built-in and program symbols are applied
directly.
%
For example, given $f$ with attribute $\rcn{t_\nil}$,
the \emph{$f$-list} with elements $\plur t n$ is the
nested application
$\eho{\eho{f}{t_1}}{\eho{\eho{f}{\ldots}}{\eho{\eho{f}{t_n}}{t_\nil}}}$.
%
The built-in symbols $\eo{nil}$, $\eo{cons}$, and
$\eo{list\_concat}$ construct and combine such
$f$-lists from within programs and rule conclusions.

A second elaboration strategy targets binder
applications.  Eunoia provides a built-in
\emph{heterogeneous list} type $\eo{List}$, with
constructors $\eo{List::nil}$ and $\eo{List::cons}$,
whose entries may have different types, and a built-in
symbol $\eo{var}$ that constructs variable-like terms
$\eapp{\eo{var}}{\estr{x}\ t}$.
%
The attribute $\bd{t_\cons}$ specifies that a binder
application $\eapp{s}{\epar{\vec ρ}\ t}$ is elaborated
so that each parameter $\epar{x_i\ t_i}$ in $\vec ρ$
becomes an $\eo{var}$ term, and these terms are
gathered with the list constructor $t_\cons$ (typically
$\eo{List::cons}$, with terminator $\eo{List::nil}$).

The $\mbf{elab}$ operator below captures both
strategies uniformly.  It is parameterized by a
signature $Σ$ and a local parameter context $\vec ρ$;
we write $\msf{progs}(Σ)$ for the set of symbols
introduced by $\kw{program}$ commands in $Σ$.

\newcommand{\body}{\text{body}}
\newcommand{\vars}{\text{vars}}
\begin{definition}\label{def:elab}
	For any signature $Σ$ and parameter context $\vec ρ$,
	let $\elab{Σ,\vec ρ}$ be the least operator such that:
	\begin{enumerate}
		\item For any symbol $s$, $\elab{Σ,\vec ρ}[s] = s$.
		      %
		\item For a binder application
		      $\eapp{s}{\epar{\vec v}\ t}$ where
		      $\prn{s ∷ \bd t_\cons} ∈ Σ$, each parameter
		      $\epar{x_i\ t_i} ∈ \vec v$ is reified as
		      $X_i ≔ \epar{\eo{var}\ \estr{x_i}\ t_i}$,
		      and the $X_i$ are gathered with the list
		      constructor $t_\cons$:
		      $$ \elab{Σ,\vec ρ}[\eapp{s}{\epar{\vec v}\ t}]
			      =
			      \eapp{s}{t_\vars\ t_\body}
			      \quad\text{where}\quad
			      \begin{array}{l@{\ {≔}\ }l}
				      t_\vars &
				      \elab{Σ, \vec ρ}[\eapp{t_\cons}{\plur X n}]
				      \\
				      t_\body &
				      \elab{Σ, \vec ρ}
				      [\mbf{subst}_{\prn{\vec x, \vec X}}[t]]
			      \end{array}
		      $$

		      %
		\item For an application $\eapp s {\vec t}$,
		      let $t_i' ≔ \elab{Σ, \vec ρ}[t_i]$ in:
		      \begin{enumerate}
			      \item If $s$ is a built-in or
			            $s ∈ \msf{progs}(Σ)$, then
			            $\elab{Σ, \vec ρ}[\eapp{s}{\vec t}] =
				            \eapp{s}{\vec t'}$.
			      \item If $s ∈ \vec ρ$, then
			            $\elab{Σ, \vec ρ}[\eapp{s}{\vec t}] =
				            \eho{\ldots\eho{\eho{s}{{t_1'}}}{\ldots}}{{t_n'}}$.

			      \item If $\prn{s ∷ α} ∈ Σ$ for an
			            associative $α$, the result is
			            one of:
			            $$
				            \begin{scases}
					            \foldr(G, t_\nil, \vec t')
					            & \tsm{$\prn{α = \rcn{t_\nil}}$}
					            \\
					            \foldl(G, t_\nil, \vec t')
					            & \tsm{$\prn{α = \lcn{t_\nil}}$}
					            \\
					            \foldr(G, t'_n, t'_1\ldots t'_{n-1})
					            & \tsm{$\prn{α = \rc}$}
					            \\
					            \foldl(G, t'_1, t'_2\ldots t'_{n})
					            & \tsm{$\prn{α = \lc}$}
				            \end{scases}
			            $$
			            where the combining function
			            $G(x,y)$ is $\econcat{s}{x}{y}$ if
			            $\prn{x ∷ \attr{list}} ∈ \vec ρ$
			            and $\eho{\eho{s}{x}}{y}$
			            otherwise.

			      \item If $\prn{s ∷ α} ∈ Σ$ for a
			            non-associative $α$, then
			            $\elab{Σ,\vec ρ}[\eapp{s}{\vec t}] =
				            \elab{Σ,\vec ρ}[\eapp{t_\cons}{\vec u}]$
			            where:
			            $$
				            \vec u ≔
				            \begin{scases}
					            \eapp{s}{t_1 t_2}\,\eapp{s}{t_2 t_3}
					            \ldots\,\eapp{s}{t_{n-1} t_n}
					            &
					            \tsm{$\prn{α = \ch {t_\cons}}$}
					            \\
					            \eapp{s}{t_i t_j} \text{ for all } i < j
					            &
					            \tsm{$\prn{α = \pw {t_\cons}}$}
					            \\
					            \vec t
					            &
					            \tsm{$\prn{α = \al {t_\cons}}$}
				            \end{scases}
			            $$

		      \end{enumerate}
	\end{enumerate}
\end{definition}

\begin{example}\label{ex:elab}
	The constants $\ev{or}$, $\ev{=}$, and $\ev{forall}$
	carry attributes that drive each of the elaboration
	cases:
	%
	\begin{eobox}
		\lstinputlisting[style=eo, aboveskip=0pt, belowskip=0pt]{listings/decls-or-eq.eo}%
		\listsep
		\lstinputlisting[style=eo, aboveskip=0pt, belowskip=0pt]{listings/decl-forall.eo}%
	\end{eobox}
	%
	\noindent
	Since $\ev{or}$ has attribute $\rcn\false$, item~3c
	applies with $G(x,y) = \eho{\eho{\ev{or}}{x}}{y}$ and
	the right fold gives:
	$$ \elab{Σ, \vec ρ}[\eapp{\ev{or}}{\ev x\ \ev y\ \ev z}]
		=
		\eho{\eho{\ev{or}}{\ev x}}{
			\eho{\eho{\ev{or}}{\ev y}}{
				\eho{\eho{\ev{or}}{\ev z}}{\false}}}
	$$
	%
	Since $\ev{=}$ has attribute $\ch{\ev{and}}$, item~3d
	first chains the equalities and then recurses
	through $\ev{and}$'s $\rcn\false$ attribute:
	$$
		\elab{Σ, \vec ρ}[\eapp{\ev{=}}{\ev a\ \ev b\ \ev c\ \ev d}]
		=
		\elab{Σ, \vec ρ}[\eapp{\ev{and}}{
				\eapp{\ev{=}}{\ev a\ \ev b}\
				\eapp{\ev{=}}{\ev b\ \ev c}\
				\eapp{\ev{=}}{\ev c\ \ev d}}]
	$$
	%
	Since $\ev{forall}$ has attribute
	$\bd{\eo{List::cons}}$, item~2 reifies each parameter
	as an $\eo{var}$ term inside an $\eo{List}$ and
	elaborates the body under the substitution
	$\sigma ≔ \prn{\ev{x} ↦ X}$, where
	$X ≔ \epar{\eo{var}\ \estr{x}\ \ev{Int}}$:
	$$
		\begin{array}{r@{\ {=}\ }l}
			\elab{Σ, \vec ρ}[\eapp{\ev{forall}}{\epar{\epar{\ev{x}\ \ev{Int}}}\ \ev{F}}]
			&
			\eapp{\ev{forall}}{t_\vars\ t_\body} \\[1mm]
			t_\vars &
			\eapp{\eo{List::cons}}{X\ \eo{List::nil}} \\[1mm]
			t_\body &
			\elab{Σ, \vec ρ}[σ\ev{F}]
		\end{array}
	$$
\end{example}

\end{document}
